%%=============================================================================
%% Samenvatting
%%=============================================================================

% TODO: De "abstract" of samenvatting is een kernachtige (~ 1 blz. voor een
% thesis) synthese van het document.
%
% Een goede abstract biedt een kernachtig antwoord op volgende vragen:
%
% 1. Waarover gaat de bachelorproef?
% 2. Waarom heb je er over geschreven?
% 3. Hoe heb je het onderzoek uitgevoerd?
% 4. Wat waren de resultaten? Wat blijkt uit je onderzoek?
% 5. Wat betekenen je resultaten? Wat is de relevantie voor het werkveld?
%
% Daarom bestaat een abstract uit volgende componenten:
%
% - inleiding + kaderen thema
% - probleemstelling
% - (centrale) onderzoeksvraag
% - onderzoeksdoelstelling
% - methodologie
% - resultaten (beperk tot de belangrijkste, relevant voor de onderzoeksvraag)
% - conclusies, aanbevelingen, beperkingen
%
% LET OP! Een samenvatting is GEEN voorwoord!

%%---------- Nederlandse samenvatting -----------------------------------------
%
% TODO: Als je je bachelorproef in het Engels schrijft, moet je eerst een
% Nederlandse samenvatting invoegen. Haal daarvoor onderstaande code uit
% commentaar.
% Wie zijn bachelorproef in het Nederlands schrijft, kan dit negeren, de inhoud
% wordt niet in het document ingevoegd.

\IfLanguageName{english}{%
\selectlanguage{dutch}
\chapter*{Samenvatting}
\selectlanguage{english}
}{}

%%---------- Samenvatting -----------------------------------------------------
% De samenvatting in de hoofdtaal van het document

\chapter*{\IfLanguageName{dutch}{Samenvatting}{Abstract}}

De European Cyber Resilience Act (CRA) legt vanaf 2026 strengere eisen op aan softwareproducenten inzake transparantie over gebruikte componenten en kwetsbaarheidsbeheer. Dit onderzoek analyseert hoe deze regelgeving kan worden geautomatiseerd binnen CI/CD-pipelines bij Excentis, een bedrijf dat reeds Jenkins pipelines gebruikt maar geen systematische securitycontroles implementeert.

De centrale onderzoeksvraag luidt: "\textit{Hoe kan een CI/CD-pipeline geautomatiseerd worden om CRA-compliant te zijn?}"

Om deze vraag te beantwoorden werd een Proof of Concept (PoC) ontwikkeld volgens DevSecOps-principes. Een referentie-Jenkins pipeline werd opgezet met open source tools: Syft voor CycloneDX SBOM-generatie, OWASP Dependency-Check voor Software Composition Analysis (SCA), SonarQube voor Static Application Security Testing (SAST) en OWASP ZAP voor Dynamic Application Security Testing (DAST). Een quality gate blokkeert builds bij CVSS > 7.0 kwetsbaarheden.

De PoC toont aan dat deze securitycontroles technisch haalbaar zijn in CI/CD-pipelines. SBOMs worden succesvol gegenereerd, scanning detecteerd meeste van de kwetsbaarheden, en de HIGH en CRITICAL kwetsbaarheden werden geblokkeerd door de quality gate.

Deze resultaten bewijzen dat CRA-conformiteit automatiseerbaar is zonder bestaande ontwikkelprocessen te verstoren. De ontwikkelde pipeline vormt een kant-en-klare referentie-implementatie voor het Excentis-team. Beperkingen zijn de scope (enkel getest op Maven, .NET en Python applicaties), performance-overhead bij grote repositories en false positives die manuele validatie vereisten.